\textit{
The integration of inverter-based resources (IBR) into transmission systems changes the fault-current characteristics used as the basis for transmission line protection. These changes can affect distance relay or relay 21, which operates based on apparent impedance, and line differential relay or relay 87L, which operates based on current comparison at both ends of the line. This study aims to analyze the operating characteristics of relay 21 and relay 87L in an IBR-integrated transmission system, compare the effects of wind turbine type 3, wind turbine type 4, and photovoltaic resources, and evaluate the feasibility of conventional protection settings under the tested scenarios. The study was conducted through DIgSILENT PowerFactory simulations using a modified two-area Kundur system. The TL89 line between Bus 8 and Bus 9 was used as the protected line. The base case without IBR was used as the reference and was compared with systems integrated with WT Type 3, WT Type 4, and photovoltaic resources. Relay 21 settings were arranged based on the conventional system condition, while relay 87L used the default settings available in the model. The evaluation considered an internal DLG fault on TL89 and external DLG, three-phase, and SLG faults on TL79 with TL78 out of service. The observed parameters included operating status, zone and apparent impedance of relay 21, and the relationship between differential current and restraint current of relay 87L. The simulation results show that relay 21 with conventional settings still operates according to the reference for an internal DLG fault on TL89. However, for external faults on TL79 with TL78 out of service, relay 21 on the Bus 8 side repeatedly fails to satisfy the starting requirement because the measured fault current is low, so its response is categorized as a misoperation. The Bus 9 side generally remains untripped because the fault is not detected in the forward direction of the relay. In the photovoltaic scenario, the apparent impedance of the ZBG element on the Bus 9 side temporarily enters zone 1, but operation is blocked by the phase directional element. In contrast, relay 87L shows a more consistent response. For the internal fault, the differential current is located in the operating region and the relay issues a trip command, while for external faults the current points remain in the restraint region and the relay does not issue a trip command. Based on these results, the conventional setting of relay 21 is feasible only in a limited manner for the system and scenarios tested, particularly when the fault current, apparent impedance, and fault direction still support relay operation. Re-evaluation is required for external faults, out-of-service line configurations, starting-current adequacy, and directional fault detection. Relay 87L shows more consistent feasibility across the tested scenario combinations, but this conclusion remains limited to the system model, IBR type and level, fault location, network configuration, inverter control, and settings used in this study.
}

\noindent\textbf{Keywords} : inverter-based resources, relay 21, relay 87L, transmission protection, conventional settings
